\documentclass[11pt]{article}
\usepackage[margin=1in]{geometry}
\usepackage{amsmath,booktabs,hyperref}
\title{Auditing Rare-Case Coverage in LLM Agents with Iterative Example Memory}
\author{Tran Vu Bao Minh\\Independent Researcher}
\date{Working draft}

\begin{document}
\maketitle

\section{Introduction}

An agent can improve its working context by storing examples that it
solves and verifying the answers before reuse. The resulting memory is
shaped by which tasks are available, which tasks the agent attempts,
which answers pass verification, and which examples the memory policy
evicts. A low-frequency task family can disappear even when each
retained answer is correct and aggregate accuracy improves. We study this
mechanism in a controlled planning domain with fixed model weights.

We distinguish coverage, the composition of memory, from generalization,
the model's accuracy on structurally disjoint held-out tasks. We compare
random and coverage-targeted task selection at matched inference budgets.
We also vary whether the added example must be solved by the model or can
use an answer from an exact solver. The latter intervention tests whether
access to a verified answer changes retention when the agent already has
an opportunity to attempt the task.

\section{Related work}

Memory studies have evaluated selective addition and deletion and how
agents reuse experience on similar tasks
\cite{experience-following}. EvoMemBench compares memory strategies
across knowledge and execution tasks \cite{evomembench}. We focus on a
narrow setting with exact answer verification and track family-level
counts at each memory update. Our intervention changes access to a
verified answer for one task while holding the task slots and model-call
budget fixed.

Studies of recursive synthetic training show that low-frequency events can
vanish across model generations \cite{shumailov2024}. Formal verification
and diverse task proposals are important in some self-play settings
\cite{psv}. Here the model parameters are fixed and the state carried
across rounds is a list of solved examples. The question is which examples
survive and whether the surviving examples affect held-out answers.

\section{Planning tasks and memory}

Each task is a directed graph with unit-cost edges, a start, a goal, and
an optional key node. A locked edge can be traversed only after visiting
the key node. A valid answer is a shortest path or a declaration that
the goal is unreachable. Breadth-first search on the product state
$(\text{location},\text{has key})$ supplies the exact shortest distance
or proves unreachability. The verifier checks the answer format, every
path transition, and optimality. The three families are ungated,
key-detour, and unreachable. Every generated graph has six nodes and
seven edges. Solvable instances have optimal distance four. In
key-detour tasks the goal is inaccessible with the key disabled, and in
unreachable tasks ignoring locks would create a path to the goal.
We define the reference frequencies as $(0.8,0.1,0.1)$. The latter two
families are rare under that reference distribution.

Let $M_{f,t}$ be the number of examples of family $f$ after update $t$.
For each family and update, the audit records attempted examples
$A_{f,t}$, answers passing verification $V_{f,t}$, admissions $I_{f,t}$,
and evictions $E_{f,t}$. The identity
\[
 M_{f,t+1} = M_{f,t} + I_{f,t} - E_{f,t}
\]
separates lost coverage from a lack of task attempts. With full FIFO
memory, if a family receives no new admissions, its last old example is
evicted once enough later examples have been admitted to reach that
example's initial position. This can occur while the combined count of
rare examples increases, because other rare families can replace it.

\section{Experimental design}

The exploratory feasibility runs compare self-selection from a fixed
menu with random or targeted supplementation. They establish the
engineering workflow and motivate the family-level audit. These runs
were used to design the follow-up and are reported separately from it.

The follow-up is a paired $2\times2$ experiment. The allocation rule
draws a family at random from the reference distribution or targets the
least represented rare family in the current memory. The answer source
is the model or the exact solver for the one external supplement. In
every arm the model is asked to solve two common background tasks and
one supplement per round. In oracle arms, we replace the supplement's
model answer before applying the shared verification and admission rules.
All background answers come from the model. Because the targeted rule
depends on memory, the answer-source contrast captures the total policy
effect across the two rounds.

We use six paired seeds with shuffled initial memory order. Each initial
memory holds four ungated, one key-detour and one unreachable example.
Memory has capacity six and uses FIFO eviction. Each seed has an
independent holdout with three tasks per family. We reject graph
isomorphs across seeds, task pools, and the earlier feasibility runs,
preserving node roles and directed edge locks when computing graph
equivalence. Initial-memory and empty-memory controls use the same
holdout within the seed. Held-out answers are not used to update memory.

The primary mechanism contrast is the paired difference in the final
number of unreachable examples between targeted oracle and targeted
model. The primary performance contrast is the paired difference in
final rare-family macro accuracy between targeted model and random
model. We report seed-level differences and their mean. An exact
two-sided sign-flip calculation is reported under sign exchangeability
as a sensitivity analysis for the small paired sample.

\section{Results}

This working draft documents the design. Follow-up results are pending.

\bibliographystyle{plain}
\bibliography{references}
\end{document}
